\chapter{Machines}
\label{ch:machines}

The single-block devices. Everything here is placed, powered and used on its own; the assembled
structures are in \autoref{ch:multiblocks}.

\section{Powered}

\begin{description}
\item[Core Sample Drill] Takes a core sample of the chunk, telling you what mineral vein is beneath
      it. The prerequisite for an Excavator.
\item[Charging Station] Charges tools, armour and any item that holds flux.
\item[Fluid Pump] Moves fluid through pipes, and is what makes a pipe network do anything.
\item[Fluid Placer] Places fluid into the world as a source block.
\item[Garden Cloche (Belljar)] Grows a crop in a glass jar from water and fertiliser. Slow, tidy,
      and entirely automatic.
\item[Floodlight] A directional lamp, the brightest light in the mod and the dearest to run. Power
      alone will not light it: a Floodlight wants a \emph{redstone signal} to switch on. Hammer it
      to aim --- which axis swings depends on the face you strike, and sneaking reverses the
      direction --- and sneak-hammer its base to invert the redstone control so it lights when the
      signal is removed instead.
\item[Turret (Gun / Chemical)] An automated defence, configurable by owner and target type.
\item[Tesla Coil] Damages entities in a radius. Genuinely dangerous to its owner.
\item[Powered Lantern] A small lamp, cheap enough to use in quantity. It is a wire endpoint in its
      own right: run an LV wire to the lantern itself, with no connector bolted to it. Do not
      confuse it with the decorative \textbf{Lantern}, which takes no power at all.
\end{description}

\begin{iefork}
City mode caps how many light blocks a Floodlight may place and stops it re-tracing its beams on a
timer --- it re-traces only when the light switches or a neighbour changes. Floodlights are usually
the most expensive block in a city build.
\end{iefork}

\section{Unpowered}

\begin{description}
\item[Wooden Barrel] Holds fluid, and is where creosote treats wood by hand.
\item[Wooden Crate] Storage that keeps its contents when broken.
\item[Reinforced Crate] The same, lockable.
\item[Workbench] Where blueprints are used and tools are modified.
\item[Item Router] Sorts items by filter into up to six outputs.
\item[Fluid Sorter] The same for fluid.
\item[Redstone Wire Connector] Carries a redstone signal down a wire rather than along the ground.
      The \textbf{Redstone Probe Connector} is its reading half, putting the state of a redstone
      network onto the block it is attached to.
\item[Turntable] Rotates whatever block is placed against it when it is given a redstone pulse.
      Hold an Engineer's Hammer and look at it to see which side it acts on and which way it turns;
      right-click with the hammer to reorient it, and sneak while doing so to reverse the
      direction.
\item[Razor Wire] A barrier that cuts anything crossing it --- 1.5 hearts bare, halved by boots or
      leggings, and a quarter with both. Strung together it forms a fence, and because it takes an
      \emph{LV wire} of its own you can electrify a run: an energised stretch shocks everything
      standing in it for a further two, at a cost of 64\,IF a pulse.
\end{description}

\section{Fluid pipes}

Steel pipes carrying fluid between machines. They connect automatically, can be covered with sheet
metal to hide them, and need a Fluid Pump to move anything.

\begin{iefork}
Two changes here. The pipe network's flood fill was rewritten --- it used a linear scan for its
closed set and removed from the head of a list for its open set, which is quadratic on a large tank
farm --- and city mode lets a pipe hand its fluid to endpoints in order rather than simulating a
fill against every one of them and splitting in proportion. Nothing is created or destroyed; only
which of several tanks fills first changes.
\end{iefork}
